\chapter{Proofs for the Constructive Spatial RG Family}
\label{app:spatial-rg-proofs}

This appendix proves the constructions of \cref{ch:constructive-spatial-rg}.
It uses that chapter's inventories, reference measures, and parameters, and
the local notation of \cref{tab:rg-lattice-notation}.

\section{The source measure and boundary comparison}
\label{sec:rg-spatial-proof-measure}

\noindent\emph{Proof of \cref{thm:rg-spatial-source-law}.}
For fixed matter signs the coordinate change \(\psi\mapsto z\) is a
bijection preserving normalized counting, with inverse
\(\psi_e=\phi_xz_e\phi_{x+e_\iota}\). The finite raw law is therefore
independent uniform matter times its dressed law in these coordinates. The
exterior \(\eta\) is dressed data; varying an exterior matter frame while
holding \(\eta\) fixed requires changing its reconstructed raw links as
well.

For \(\Delta\subset\Lambda\), conditioning on \(z_{\Lambda\setminus\Delta}\)
cancels the edge factors outside \(\Delta\) and the plaquettes not meeting
\(\Delta\). The remaining factors are exactly \(W_\Delta\) with the
completed exterior, so the finite specification is consistent. Its one-edge
conditional plus probability is
\begin{equation}
 q_e=\frac{1+\tanh\widetilde h_e}{2},\qquad
 \widetilde h_e=h+j\sum_{p\ni e}\prod_{f\in p\setminus\{e\}}z_f,
 \qquad |\widetilde h_e|\le h+2j.
 \label{eq:rg-spatial-proof-single-row}
\end{equation}
Every conditional sign probability is at least \(\epsilon_f=(1-\rho)/2>0\).
Averaging over unrevealed variables preserves this bound, so a specified
assignment on \(A\) has probability at least \(\epsilon_f^{|A|}\).

Let \(n(e,f)\) count the plaquettes containing both of the distinct edges
\(e,f\). Each edge lies in two plaquettes with three other edges in each,
whence \(\sum_fn(e,f)=6\). Changing one plaquette contribution changes the
conditional mean by
\begin{equation}
 \frac{|\tanh(s+j)-\tanh(s-j)|}{2}
 =\frac{\sinh(2j)}{\cosh(2s)+\cosh(2j)}\le t.
 \label{eq:rg-spatial-proof-influence}
\end{equation}
Telescoping such changes gives the influence matrix \(D_{ef}=t\cdot n(e,f)\),
with row sums at most \(\alpha=6t<1\).

A random-coordinate heat bath on the finite free set preserves its Gibbs
law by conditional factorization. For \(n=|\Lambda|\), updating each
coordinate in a prescribed order to any target has probability at least
\((\epsilon_f/n)^n\), so every entry of the \(n\)-step transition matrix
has that positive lower bound; removing the common row mass gives a strict
total-variation contraction and a unique stationary law. Couple two such
chains by the same chosen edge and the same uniform number in their
conditional Bernoulli rows. Ces\`aro averages of the finite pair chain have
a convergent subsequence, whose limit is stationary because one application
changes the average by a difference divided by its length; its marginals
are the two Gibbs laws. If \(d_e\) is its probability of disagreement at
\(e\), stationarity cancels the unchanged update probability and gives
\begin{equation}
 d_e\le\sum_{f\in\Lambda}D_{ef}d_f+\beta_e,
 \qquad \beta_e=\sum_{f\notin\Lambda}D_{ef}
                           \mathbf1_{\{\eta_f\ne\eta'_f\}}.
 \label{eq:rg-spatial-proof-stationary-comparison}
\end{equation}
Iteration removes the remainder \(D_\Lambda^nd\), which is bounded by
\(\alpha^n\). The surviving sum runs over paths making their first exit
from \(\Lambda\) at a differing exterior edge. Dividing each transition
weight by \(\alpha\), and adding a cemetery state for missing row mass,
turns their weights into disjoint first-exit probabilities with sum at most
one. Every exit path from \(e\) has at least \(\operatorname{dist}_\square(e,\Lambda^c)\) steps, so
\(d_e\le\alpha^{\operatorname{dist}_\square(e,\Lambda^c)}\), and a union bound over target
disagreements proves \cref{eq:rg-spatial-boundary-comparison}.

A larger finite-volume restriction is a mixture of the smaller
specifications, by finite consistency, so integrating the bound over two
boundary mixtures compares volumes containing a common interior. On growing
boxes each fixed cylinder is therefore Cauchy, independently of all
exterior choices, and its limit is normalized, consistent, and has every
atom at least \(\epsilon_f^{|A|}\). To construct the measure, enumerate the
countable edges, split \([0,1)\) into half-open intervals with the
probabilities of the first edge, and split each interval into the joint
probabilities of the next edge; consistency makes every subdivision exact.
Each point specifies a unique infinite sign sequence whose cylinder
preimages are finite unions of intervals, so the pushforward of Lebesgue
measure is a countably additive law with these cylinders. Cylinders form a
\(\pi\)-system, and the \(\pi\)--\(\lambda\) theorem makes the extension
unique.

For a cylinder function \(f\), the finite update \(\varsigma_\Lambda f\)
remains a cylinder function because the interaction boundary is finite.
Passing finite compatibility to the cylinder limit gives
\(\mu f=\mu\varsigma_\Lambda f\), and cylinder determination extends the
identity to bounded measurable functions. For bounded exterior-measurable
\(w\), \(\varsigma_\Lambda(wf)=w\varsigma_\Lambda f\), so this is the
conditional DLR identity. Any two DLR laws have target marginals that are
mixtures of the same finite specifications; the boundary bound tending to
zero makes all their cylinders, hence the laws, equal. A translate of the
unique law is a law for the same specification and therefore agrees with
it. The conditional law in an ordinary box embedded in a full torus has the
same factors and a random dressed exterior, so the same comparison proves
the torus cylinder limit while retaining all winding sectors.

For the raw lift, append independent uniform matter signs and apply the
inverse coordinate map. A finite raw window with free matter vertices and
all their incident links converts to independent free frames and the
correct dressed conditional; such windows are cofinal, and conditioning
further proves the raw DLR equations for every finite window. Conversely, a
local gauge flip of one matter vertex and its incident links preserves each
raw conditional density and reference, so any raw Gibbs law is invariant
under every finite product of these flips. In dressed coordinates these
flip arbitrary finite sets of matter signs without changing \(z\).
Averaging over the flips shows that every assignment of the matter signs of
\(k\) given sites has probability \(2^{-k}\), independently of every dressed
cylinder. The dressed marginal has the specification already identified,
hence equals \(\mu_{b,t}\), which proves uniqueness of the raw lift.
\hfill\(\square\)

\section{Constrained parity laws uniformly in blocking depth}
\label{sec:rg-spatial-proof-constrained}

\noindent\emph{Proof of \cref{thm:rg-spatial-constrained-specification}.}
Partition the fine edges into the retained \(B\)-paths and unused singleton
edges. No plaquette contains two edges of one retained path. A coarse
anchor cell contains two retained paths and \(2B(B-1)\) unused singleton
blocks, hence \(2B^2\) fine edges. A pattern prescribes each retained block
to be unconstrained or to have one specified parity; singleton blocks
remain unconstrained. Use the product of uniform probability references on
these supports, multiply all edge and touching plaquette factors, and
normalize; these are positive finite specifications on their supports. A
block row with at most \(B\) edges has every allowed atom at least
\begin{equation}
 \xi_B=2^{-B}\left(\frac{1-b}{1+b}\right)^B
                   \left(\frac{1-t}{1+t}\right)^{2B}>0.
 \label{eq:rg-spatial-proof-block-positivity}
\end{equation}

Two product sign laws obtained from one positive product law by
conditioning on opposite total parities admit a coupling at Hamming
distance exactly one. For one bit this is immediate. For \(n\) bits, couple
the last-bit marginals by any coupling. If those bits differ, the required
prefix parities agree and the prefixes are coupled identically; if they
agree, the induction coupling is applied to the opposite prefix parities.
In either branch the prefix marginal on each side is its product law
conditioned on the required parity, independently of the other branch,
which proves the induction including both marginals.

Changing a single coordinate field from \(+j\) to \(-j\) in a product row
or fixed-parity product row changes that coordinate's marginal by at most
\(t\), by \cref{eq:rg-spatial-proof-influence} with the other-coordinate
masses absorbed into \(s\). Couple it maximally. If it differs in a
fixed-parity row, the remaining prefix parities are opposite, so the
preceding coupling adds exactly one further discrepancy. The expected
Hamming cost is at most \(2t\), and a one-bit fixed fiber has zero cost.
Successive tilts are coupled by composing their conditional coupling
kernels, and the endpoint Hamming distance is bounded by the sum of
successive distances, which is the required triangle inequality.

For a target block \(i\) and an exterior fine edge \(e\), let \(n(i,e)\) be
the number of plaquettes meeting both. An exterior flip changes at most
\(n(i,e)\) single-coordinate fields, so the preceding coupling gives a row
cost at most \(2t\sum_en(i,e)\mathbf1_{\{z_e\ne z'_e\}}\). The column
identity, not a block-state row bound, is \(\sum_in(i,e)=6\): each of the
two plaquettes of \(e\) meets three other blocks. Incident anchors differ
by at most one.

Finite block heat baths have unique stationary laws by the common-mass
argument using \(\xi_B\), and a stationary pair coupling exists by the same
Ces\`aro construction as above. Let \(d_e\) be the fine-edge disagreement
probabilities and \(D_i=\sum_{e\in i}d_e\). Canceling the update probability
separately in each block gives
\begin{equation}
 D_i\le2t\sum_{e\text{ free}}n(i,e)d_e
      +2t\sum_{e\text{ exterior}}n(i,e)
                         \mathbf1_{\{\eta_e\ne\eta'_e\}}.
 \label{eq:rg-spatial-proof-block-hamming}
\end{equation}
For a target anchor set \(A\), multiply by
\(w_i=e^{-\tau\operatorname{dist}_\infty(\operatorname{anchor}(i),A)}\) and
sum. The column identity and \(w_i\le e^\tau w_{\operatorname{block}(e)}\)
when \(n(i,e)>0\) yield
\begin{equation}
 S\le\mathrm qS+\mathrm q\sum_{e\in\partial_f\Lambda}w_{\operatorname{block}(e)}
                          \mathbf1_{\{\eta_e\ne\eta'_e\}},\qquad
 S=\sum_{i\text{ free}}w_iD_i.
 \label{eq:rg-spatial-proof-weighted-column}
\end{equation}
Use all blocks anchored in the radius-\(R\) neighborhood of \(A\). Every
contributing exterior edge is anchored in its next shell, whose cardinality
is at most \(8|A|(R+1)\), because a square shell around one cell has exactly
\(8(R+1)\) cells, and there are \(2B^2\) fine edges per anchor. The exterior
sum is therefore at most \(16|A|B^2(R+1)e^{-\tau(R+1)}\). Dividing by
\(1-\mathrm q\) and taking a union bound over target disagreements proves
\cref{eq:rg-spatial-coarse-modulus} for these finite constrained laws.

For every fixed full pattern, the bound makes all finite-cylinder limits
Cauchy on expanding whole-block volumes. Interval subdivision constructs
their measure, finite compatibility passes to the limit and gives its DLR
law, and any two such laws have cylinders within the same vanishing
boundary bound and hence agree. The argument also compares patterns
agreeing in a large interior; different supports outside it are only
exterior forcing. No conditioning on an infinite event is used.

Finite conditioning on previously free path parities has positive
probability by \cref{eq:rg-spatial-proof-block-positivity}. Restricting its
finite conditional reference is precisely the updated support pattern, and
uniqueness identifies the resulting infinite law. In the all-unconstrained
pattern the law is the original fine Gibbs law, because its finite DLR
identities also hold on unions of blocks, so its finite parity conditionals
are these constrained laws. To define a target coarse kernel, leave its
paths unconstrained, prescribe every exterior retained parity, and take the
target marginal of the constructed constrained law.

Conditional on all other block states, the signs along a free path are
independent, with conditional means bounded by \(\rho\). Its parity mean
has magnitude at most \(\rho^B\), also after conditioning on any already
revealed disjoint parities. Revealing the two target parities successively
gives the asserted one-cell lower and upper probabilities. The boundary
comparison makes every kernel continuous and measurable, and conditioning
on a finite set of additionally revealed parities identifies the
corresponding updated pattern, so finite Bayes identities give pointwise
specification consistency at every exterior.

For an exterior cylinder event, the defining finite conditional integral
holds under \(\mu_B\). Increasing the revealed exterior set and using
uniform convergence passes the integral to the constructed kernel, and
cylinder determination extends this conditional identity to the full
exterior \(\sigma\)-algebra. Thus \(\mu_B\) is DLR for this
everywhere-defined specification. Any two coarse DLR laws can be mixed over
increasingly remote exteriors, and the same uniform comparison makes their
target cylinders equal, which proves uniqueness. Translation covariance
follows by transporting factors and patterns and invoking constrained
uniqueness. At \(t=0\), independent product-fiber laws give a separate
exact control with coarse edge mean \(b^B\), and no formula dividing by
\(\mathrm q\) is needed. \hfill\(\square\)

\section{The complete canonical interaction}
\label{sec:rg-spatial-proof-potential}

\noindent\emph{Proof of \cref{thm:rg-spatial-generated-interaction}.}
The one-cell bounds give log-odds amplitude
\(4\operatorname{atanh}(\rho^B)\le\mathrm c_2\rho^B\). Atom differences are
bounded by total variation, and \(|\log s-\log s'|\le|s-s'|/p_{\min}\) for
\(s,s'\ge p_{\min}\), so reference-state log odds have spatial variation at
most \(\mathrm c_3B^2(R+1)e^{-\tau(R+1)}\). The canonical supports have
diameter at most \(2r\), and at most \((2r+1)^2\) translates contain the
origin. Their maximal lexicographic cell uniquely specifies the anchor, so
no assignments collide. Consequently
\begin{equation}
 \|\mathscr U_{B,A_{x,r}}\|_\infty
 \le\min\{2\mathrm c_2\rho^B,\mathrm c_3B^2re^{-\tau r}\},\qquad r\ge1.
 \label{eq:rg-spatial-proof-shell-minimum}
\end{equation}
The inequality \(\min(s,s')\le\sqrt{ss'}\), the bound \(\sqrt r\le r\), and
\(\lambda_v=\tau/8\) reduce the weighted shell sum to
\(\sqrt{2\mathrm c_2\mathrm c_3}B\rho^{B/2}\sum_{r\ge1}r(2r+1)^2z_v^r\).
Differentiating the geometric series three times gives exactly
\(z_v(9+14z_v+z_v^2)/(1-z_v)^4\), which proves the norm bound; its explicit
tail after \(N\) shells replaces that series with
\(\sum_{r>N}r(2r+1)^2z_v^r\). Since
\(\operatorname{diam}A\le e^{\lambda_v\operatorname{diam}A}/\lambda_v\), it
also proves the first moment. Both terms in \(\mathcal V_B\) tend to zero;
for an exponential bound, \(B\rho^{B/4}\le\rho^{1/4}/(1-\rho^{1/4})^2\).

For a configuration differing from \(*\) on finitely many cells, change
those cells in increasing lexicographic order. Positive finite
specification consistency expresses its log probability ratio against the
vacuum as the sum of the corresponding full-past log odds, and the
absolutely convergent canonical shells telescope to exactly those odds.
Terms anchored outside the changed set vanish, since their anchor stays at
\(*\). The interaction and the specification therefore give identical
vacuum-exterior ratios. Comparing two finite-support configurations and
canceling their common exterior gives equality for every finite-support
exterior. Approximate an arbitrary exterior by its values in larger boxes
and vacuum outside: quasilocality gives uniform convergence on the
specification side, and absolute summability of all terms touching the
finite changed set gives uniform convergence of interaction energy
differences. Positive normalization then gives equality for every exterior.
No infinite total Hamiltonian is used as an unnormalized density.

A continuous specification of a full-support law is unique everywhere: a
difference at one exterior would persist on a cylinder of positive
probability, contradicting almost-sure equality of conditional versions.
The canonical interaction therefore depends only on the law. Path
composition and this uniqueness give the stated action on its orbit, even
if two block sizes happen to give the same interaction. At the fair product
law all odds are zero, so the interaction is zero and fixed. At \(t=0\) the
singleton interaction and its normalized-counting pressure are
\begin{equation}
 \mathscr U_{B,\{x\}}(\sigma)=-\sum_{\iota=1}^2
       \log\frac{1+b^B\sigma_\iota}{1+b^B},\qquad
 \mathfrak p(\mathscr U_B)=-2\log(1+b^B),
 \label{eq:rg-spatial-proof-product-vacuum-scalar}
\end{equation}
where \(\sigma=(\sigma_1,\sigma_2)\in\{-1,1\}^2\) is the cell state. This finite scalar is
not zero despite vacuum normalization. \hfill\(\square\)

\section{The open-ball coefficient map and its law}
\label{sec:rg-spatial-proof-operator}

\noindent\emph{Proof of \cref{thm:rg-spatial-local-operator}.}
The coefficient space is complete: a norm-Cauchy sequence converges in
every anchored weighted \(\ell^1\) space, its coordinate limits agree on
overlaps, and the uniform Cauchy bound passes to the supremum over anchors.
Translation invariance is a closed coordinate condition. Finite
translation-orbit families of arbitrarily large supports give linearly
independent members, so the positive-radius ball is infinite dimensional.

Condition the fair fine reference on all retained \(a\)-path parities. Its
blocks are independent uniform parity fibers and unused fair singletons. A
character has conditional expectation zero unless its support is a union of
whole paths: a partially occupied path is killed by flipping one included
and one excluded sign, and an occupied unused singleton by its own flip. A
surviving support \(\operatorname{lift}_a(T)\) has expectation \(\chi_T\).
For an ordered list of source occurrences \(S_1,\ldots,S_n\), repetitions
allowed, every product of moments in the joint cumulant is therefore either
zero or the character of
\(S_1\mathbin\triangle\cdots\mathbin\triangle S_n\). Denote its constant
coefficient by \(c_a(S_1,\ldots,S_n)\). The degree-\(n\) map is
\begin{equation}
 P_{a,n}(J_1,\ldots,J_n)_T
 =\frac1{n!}\sum_{S_1\triangle\cdots\triangle S_n=\operatorname{lift}_a(T)}
 c_a(S_1,\ldots,S_n)\prod_{i=1}^n(J_i)_{S_i},\qquad T\ne\varnothing.
 \label{eq:rg-spatial-proof-cumulant-map}
\end{equation}
For finite source inventories this is differentiation of the finite
conditional log partition, including its \(1/n!\) and repeated
occurrences.

Join two occurrence vertices when their supports meet a common reference
block; disconnected groups then depend on independent reference variables.
The bounded dependency-graph cumulant inequality
\citep[Sec.~9.4.2, Eq.~(47)]{feray2016modphi} gives
\(|c_a|\le2^{n-1}\operatorname{ST}(G)\), where \(\operatorname{ST}\) counts
spanning trees of this occurrence graph. Its hypotheses hold uniformly in
every parity assignment: the variables are sign characters of magnitude one,
and the independence is at the block level. Convergence and the
infinite-law identification are proved below.

Write \(s_i=|S_i|\) and \(d_i=\operatorname{diam}S_i\). For nonempty output
\(T\), disjoint paths give \(a|T|\le\sum_is_i\). Along an occurrence tree,
travel within a source at cost at most \(d_i\), and between two source
edges in one reference block at cost at most \(a-1\). The first fine edges
of two output paths realizing its diameter then give
\begin{equation}
 a\operatorname{diam}T\le\sum_id_i+(a-1)(n-1).
 \label{eq:rg-spatial-proof-tree-diameter}
\end{equation}
Fix an output edge and its first fine edge. At least one occurrence
contains that fine edge, so sum over \(n\) possible roots. Pay \(|T|\) by
\(a^{-1}\sum_is_i\), summing over \(n\) marked occurrences. A child support
touching a parent reference block contains an edge in the union of blocks
met by the parent, at most \(as_i\) edges. Expanding this witness condition
and summing from the leaves costs \(a\) per tree edge and one parent
support-size factor per child.

For an occurrence with \(k_i\) children, the weighted anchored sum is
bounded by its input norm times \(k_i!\delta_a^{-k_i}\). Indeed the
unmarked residual factor is \(s_i^{k_i-1}e^{-\delta_as_i}\) and the marked
one is \(s_i^{k_i}e^{-\delta_as_i}\); both are at most
\(k_i!\delta_a^{-k_i}\), using \((\delta_as_i)^{k_i}/k_i!\le
e^{\delta_as_i}\) and \(s_i\ge1\). The unused spatial weight only decreases
the bound. Summing \(\prod_ik_i!\) over labeled trees with fixed root gives
\((n-1)!\operatorname{Cat}_{n-1}\): ordering the children counts plane
rooted trees, and the other \(n-1\) labels are free. Balanced-parenthesis
encoding gives \(\operatorname{Cat}_{n-1}\le4^{n-1}\). Combining the
cumulant tree inequality, roots, mark, \(1/n!\), and
\cref{eq:rg-spatial-proof-tree-diameter} proves
\begin{equation}
 \|P_{a,n}(J_1,\ldots,J_n)\|
 \le\frac na\mathcal K_a^{n-1}\prod_{i=1}^n\|J_i\|.
 \label{eq:rg-spatial-proof-multilinear-majorant}
\end{equation}
All rearrangements were of positive majorants, so they are justified also
for infinite source inventories.

Define \(\mathscr R_aJ=\sum_{n\ge1}P_{a,n}(J,\ldots,J)\) for
\(\mathcal K_a\|J\|<1\). The majorant gives locally uniform convergence in
the Banach norm and locally uniform convergence of
\(D\mathscr R_a(J)H=\sum_{n\ge1}nP_{a,n}(H,J,\ldots,J)\) in operator norm.
The multilinear increment expansion leaves a uniformly summable remainder
with at least two increments, so this is the Fr\'echet derivative, and its
locally uniform series is continuous. The first-order survival rule gives
\(L_a\). The lift is injective, has size \(a|T|\), contains the first edge
of every output path, and has diameter at least \(a\operatorname{diam}T\),
so comparing anchored weights gives \(\|L_a\|\le\mathfrak c_a\). For
\(\omega=\mathcal K_a\|J\|\le1/2\), \(\sum_{n\ge2}n\omega^{n-2}\le6\) and
\(\sum_{n\ge2}n^2\omega^{n-2}\le22\), by differentiating the geometric
series. These prove the quadratic bound with \(6\mathcal K_a/a\) and the
derivative bound with \(22\mathcal K_a/a=2M_a\), hence the stated bounds.
On the ball,
\(\|\mathscr R_aJ\|\le(\mathfrak c_a+M_a\|J\|)\|J\|\le\frac{1+\mathfrak c_a}2\|J\|\),
because \(\varepsilon\le(1-\mathfrak c_a)/(2M_a)\); a fixed point \(J\) in the ball therefore satisfies
\(\|J\|\le\frac{1+\mathfrak c_a}2\|J\|\) and is zero.

It remains to identify a normalized law. For any real input in the ball,
its one-edge half-log-odds are \(\sum_{S\ni e}J_S\chi_{S\setminus\{e\}}\),
absolutely and uniformly convergent. Flipping \(f\) changes the Bernoulli
row by at most \(\sum_{S\ni e,f}|J_S|\), so the influence row sums are at
most \(\|J\|<1\). Finite heat-bath coupling gives the resolvent comparison
\(d\le\sum_{m\ge0}D^m\beta\). For a fixed target, the contribution from a
remote boundary tends to zero: each fixed power has summable rows, and the
remaining powers have a common geometric tail. Finite-cylinder
subsequences exist by finite-state compactness, and uniform convergence of
finite local conditional energies passes their DLR equations to the limit.
The comparison then makes the limit unique and makes the entire exhaustion
converge.

Take finite edge boxes \(E_m\) and source truncations
\(J^{[m]}_S=J_S\mathbf1_{\{S\subset E_m\}}\). Each is a finite inventory of
no larger norm; no global norm convergence of these inhomogeneous
truncations is asserted. Every fixed source tuple is eventually present, so
the rooted majorant gives, for every fixed output edge, weighted incident
convergence of \(\mathscr R_aJ^{[m]}\) to \(\mathscr R_aJ\). In the finite
inventory the convergent scalar and nonempty Taylor series exponentiate to
the conditional moment generating function near zero. The identity of the
convergent analytic series extends over its real convergence interval, on
which that function is strictly positive, so its real logarithm has
precisely these coefficients throughout the ball. The scalar cancels only
when the finite density is normalized.

The truncated fine law is the free-boundary Gibbs law on \(E_m\), with fair
signs outside. A path crossing its boundary retains those outside fair
signs; incomplete characters are canceled rather than discarded by an
altered channel. Its path push has exactly the finite coefficient expansion
above. Incident convergence makes every finite coarse energy difference,
hence every normalized conditional row, converge uniformly in the exterior.
The fine laws converge locally to their unique law, and the path map
preserves convergence of cylinder observables. Passing the finite coarse
DLR integral to this limit proves that the path-pushed law is Gibbs for
\(\mathscr R_aJ\). Its norm lies in the ball, so it is the unique coarse
Gibbs law by the same comparison.

Every prescribed parity exterior has a continuous conditional version. In
the whole-reference-block representation, changing one outside block gives
an influence row sum bounded by \(2a\|J\|<1\), since a block has at most
\(a\) fine edges. This estimate holds on every fixed-parity or
unconstrained support, with all crossing sources present. Its summable
resolvent constructs unique laws for every pattern and gives local
continuity when remote patterns change. Free retained parity probabilities
are at least \(e^{-2a\|J\|}/2>0\), and finite parity conditioning
corresponds exactly to these product-fiber supports. Testing the uniform
limit of its finite conditional rows against exterior cylinders identifies
it as a conditional version of the path-pushed law. The coefficient Gibbs
kernel is also continuous, and since the law has full cylinder support the
two versions agree everywhere.

Finally, integrating the conditional half-log-odds at \(e\) against the
fair exterior character \(\chi_{S\setminus\{e\}}\) recovers \(J_S\) by
orthogonality and uniform convergence, which proves uniqueness of every
nonempty coefficient. Path composition, unique laws, and this extraction
give \(\mathscr R_2\mathscr R_3=\mathscr R_3\mathscr R_2=\mathscr R_6\), and
the same reasoning applies to any finite word with its directly defined
path channel. The original field and plaquette source has norm
\(e^\kappa h+8e^{4\kappa+\lambda}j\), so the ball contains a nonempty region
with both couplings positive. \hfill\(\square\)

\section{Complete finite Taylor calculations and their errors}
\label{sec:rg-spatial-proof-finite-amplitude}

\noindent\emph{Proof of \cref{thm:rg-spatial-finite-amplitude}.}
Represent each nonempty support once modulo integer translation, by moving
its lexicographically first edge basepoint to zero and retaining its axis.
If its shape is \(S\), with size \(|S|\) and diameter \(\operatorname{diam}
S\), and it contains \(n_\iota(S)\) axis-\(\iota\) edges, exactly
\(n_\iota(S)\) distinct translates contain a specified axis-\(\iota\) edge.
A nonempty finite set has no nonzero translation period, so a finite
library's norm is
\begin{equation}
 \|J\|=\max_{\iota=1,2}\sum_S n_\iota(S)|S|
                  e^{\kappa|S|+\lambda\operatorname{diam}S}|J_S|.
 \label{eq:rg-spatial-proof-finite-orbit-norm}
\end{equation}
Duplicates are consolidated by exact addition before this formula is used.
Only translations are identified; rotations and reflections are not.

At order \(n\), let \(d_{\max}\) and \(s_{\max}\) be the largest input
shape diameter and size. A connected source-occurrence list contributing an
output with first edge at the origin stays within fine basepoint radius
\(R_n=nd_{\max}+(n-1)(a-1)\), by the occurrence-tree path argument, so
every contributing literal instance belongs to the finite catalog of all
input shapes translated by \([-R_n,R_n]^2\cap\mathbb Z^2\). Reference-block
ownership uses Euclidean division: a horizontal edge is retained exactly
when its transverse coordinate is divisible by \(a\), and its owner uses
the floor of its parallel coordinate divided by \(a\); vertical edges have
the rotated rule, and negative coordinates follow the same floor rule.
Unused singleton owners have a distinct type. A support is a union of
complete paths exactly when every owner is retained and has all \(a\)
distinct members.

For a set \(I\) of occurrence labels, let \(m(I)\) be one if their
symmetric-difference support is a union of complete paths, including the
empty support, and zero otherwise; at all positive parities this is its
conditional moment. With \(r\) the first occurrence label in a nonempty
\(I\), the cumulant obeys
\begin{equation}
 c(I)=m(I)-\sum_{\substack{D\subsetneq I\\r\in D}}
                           c(D)m(I\setminus D).
 \label{eq:rg-spatial-proof-cumulant-recurrence}
\end{equation}
The partition block containing \(r\) uniquely decomposes the
moment--cumulant formula, and summing all other blocks gives the second
moment factor; this proves the recurrence and its singleton base.
Occurrence labels remain distinct even when their supports repeat.
Computing all subsets takes finitely many exact integer operations, and
fewer than \(3^n\) subset-pair contributions suffice.

For each order and each output first axis, enumerate every size-\(n\)
multiset of literal catalog instances once, and keep connected occurrence
graphs with nonempty complete-path support whose output is already
canonically anchored on that axis. A multiset with multiplicities
\(m_k\) contributes exactly
\(c(I)\prod_kJ_k^{m_k}/\prod_km_k!\), with \(k\) running over the instances: its
\(n!/\prod_km_k!\) ordered occurrences cancel the \(n!\) in
\cref{eq:rg-spatial-proof-cumulant-map}. Root choices and spanning trees are
bounds, not extra multiplicities. Disconnected occurrences have zero
cumulant by block independence, and the radius argument loses no connected
term. Output canonicalization selects the already canonical output rather
than accumulating all translated copies. Consolidation over all accepted
terms gives the coefficient at every support through the specified order.
Surviving supports satisfy \(a|T|\le ns_{\max}\) and
\(a\operatorname{diam}T\le nd_{\max}+(a-1)(n-1)\), so the output is finite,
and every finite catalog, multiset list, recurrence, and rational operation
terminates at fixed order.

Summing \cref{eq:rg-spatial-proof-multilinear-majorant} over orders above
\(\bar n\) gives \cref{eq:rg-spatial-taylor-tail} by differentiating the
geometric tail \(\omega^{\bar n+1}/(1-\omega)\). Its positive-series
representation makes it monotone in the norm and in \(\mathcal K_a\), so
upper bounds for both give a valid bound despite the minus sign in the
closed numerator. It tends to zero for \(\omega<1\), and the zero input
gives the zero table and zero tail.

The transcendental constants admit convergent rational enclosures. For a
nonnegative rational \(\xi\), put \(w_0=1\), \(w_{k+1}=w_k\xi/(k+1)\), and
\(W_m=\sum_{k=0}^mw_k\). Once \(v_m=\xi/(m+2)<1\), positivity and
decreasing successive term ratios give
\begin{equation}
 W_m\le e^\xi\le W_m+\frac{w_{m+1}}{1-v_m}.
 \label{eq:rg-spatial-proof-exp-enclosure}
\end{equation}
The width tends to zero, and reciprocal endpoints enclose \(e^{-\xi}\).
Applying these termwise to \cref{eq:rg-spatial-proof-finite-orbit-norm} and
\cref{eq:rg-spatial-local-radius} gives convergent upper norm and lower
radius bounds, so every strict admission or positive tolerance gap is
decided at some finite precision, and the strict image margin and vanishing
Taylor tail eventually admit the approximate image as well.

For one radius \(r<\varepsilon\), the derivative estimate along a segment
gives the Lipschitz constant \(\operatorname{Lip}_a(r)<1\), while the value
estimate gives the radial factor \(q_a(r)=\mathfrak c_a+M_ar<1\). The
triangle inequality separates one step's Taylor error from its input error
and gives the stated recurrence. Choosing every tail below both a
prescribed small fraction of the total tolerance and \((1-q_a(r))r/2\) keeps
all approximate states in the ball and achieves the requested error for a
finite word. Two differently truncated direct and staged tables need not be
equal; their difference is bounded by the sum of their errors relative to
the same composed operator. The empty character is excluded only from the
coefficient table: an unused fine field has a nontrivial \(\cosh\)
normalizer although its retained interaction is zero, so no evidence value
can be inferred from that table. \hfill\(\square\)

\section{A constructive basin-entry theorem in the prescribed norm}
\label{sec:rg-spatial-proof-basin}

\noindent\emph{Proof of \cref{thm:rg-spatial-basin-entrance}.}
The proof uses the compact box and retains the prescribed \(\kappa,\lambda\).
For \(B\ge3\), conditionally integrate all but the first and last edges of
each retained path. No plaquette meets two eliminated edges, and the other
three edges of such a plaquette lie outside every retained path, so these
interiors are conditionally independent. For a horizontal path at the
origin the conditional parity mean is
\begin{equation}
 \bar z_p=z_{(0,0,1)}z_{(B-1,0,1)}
        \prod_{s=1}^{B-2}f(\chi^+_s,\chi^-_s),\qquad
 f(\chi,\chi')=\tanh(h+j(\chi+\chi')),
 \label{eq:rg-spatial-proof-basin-strip}
\end{equation}
where \(\chi^\pm_s\) are the three-edge characters of the plaquettes above
and below the edge \((s,0,1)\). Its strip consists of its endpoints, the
horizontal edges at \((s,\pm1)\), \(1\le s\le B-2\), and the vertical edges
at \((s,0),(s,-1)\), \(1\le s\le B-1\). It has \(4B-4\) edges and diameter
at most \(B\). These strips and their rotations are pairwise disjoint: for
horizontal edges, horizontal-strip interiors have first-coordinate residues
\(1,\ldots,B-2\), whereas vertical-strip interiors have residues \(0,B-1\),
and the transverse residues identify the unique adjacent grid line for
\(B\ge3\). Retained endpoints belong only to their own path, and rotation
proves the other orientation. This also proves conditional independence of
all eliminated interiors.

Set \(t_2=\tanh2j=2t/(1+t^2)\). The four Walsh coefficients of \(f\) are
\(\hat f_\varnothing,\hat f_1,\hat f_1,\hat f_{12}\), with
\begin{equation}
 \hat f_\varnothing-\hat f_{12}=b,\qquad
 2\hat f_1=\frac{t_2(1-b^2)}{1-b^2t_2^2},\qquad
 \hat f_{12}=-\frac{bt_2^2(1-b^2)}{2(1-b^2t_2^2)}.
 \label{eq:rg-spatial-proof-strip-coefficients}
\end{equation}
Here \(\hat f_\varnothing,\hat f_1\ge0\) and \(\hat f_{12}\le0\), so the
unweighted coefficient sum is \(R(b,t_2)=b+t_2(1-b^2)/(1-b^2t_2^2)\). It
increases in \(t_2\), and \(\partial_bR\ge1-2t_2/(1-t_2^2)>0\) for
\(t_2\le24/145\), the value at \(t=1/12\). Since \(R(1,t_2)=1\), its
upper-endpoint value \(R_*\) bounds the compact box with the explicit slack
\(1-R_*\ge(13489/20449)(1-b_+)>0\). Put
\begin{equation}
 \Delta=1-R_*,\qquad \alpha=\Delta/24,\qquad
 \bar R=\frac{R_*}{1-\Delta/4}<1,\qquad
 \mathbb A_B=\frac{\bar R^{B-2}}{1-2\alpha}.
 \label{eq:rg-spatial-proof-basin-activity}
\end{equation}
The fine Walsh norm with weight \(e^{\alpha|S|}\), including its scalar, is
submultiplicative under symmetric difference. Since \(\chi,\chi'\) have
three edges each, \(\|f\|_\alpha\le e^{6\alpha}R_*\le\bar R\) and
\(\|\bar z_p\|_\alpha\le\mathbb A_B\), using \(e^x\le1/(1-x)\) for
\(0\le x<1\), including the two endpoints.

For a finite path set \(P\), the parity density relative to fair counting
is
\begin{equation}
 f_P(y)=\E_{\mu_{b,t}}\prod_{p\in P}(1+y_p\bar z_p).
 \label{eq:rg-spatial-proof-soft-parity-density}
\end{equation}
Conditional independence of the eliminated interiors proves this identity,
and averaging over \(y\) gives one. The auxiliary fine law tilted by this
product has the correct remaining-sign marginal for these observations; it
is not the hard-parity posterior on all fine variables, which matters for
its normalized conditional rows below.

For a root path \(p_0\) with fine anchor \(x_0\), let \(D_{p_0}(S,T)\) be
the maximum fine-coordinate distance from \(x_0\) of any fine-edge anchor in
\(S\), or of \(B\) times any coarse-edge anchor in \(T\); an empty maximum
is zero. Set \(\beta=\kappa+1\) and \(\lambda_B=2\lambda/B\), and complete
the joint fine and coarse Walsh algebra in
\begin{equation}
 \|F\|_{p_0}=\sum_{S,T}|F_{S,T}|
     e^{\alpha|S|+\beta|T|+\lambda_BD_{p_0}(S,T)}.
 \label{eq:rg-spatial-proof-rooted-algebra}
\end{equation}
Fine signs and coarse labels are distinct variables. Cardinalities and
rooted radii are subadditive under multiplication, which proves the Banach
algebra inequality. For a local update at a fine edge use the same norm
centered there. Removing that edge lowers the fine weight by \(e^{-\alpha}\),
and replacing it by a local polynomial increases the rooted radius by at
most that polynomial's local radius.

Put \(\bar p=(1+\bar R)/2\), and increase \(B\) until \(2\lambda/B<1\) and
\(\bar R/(1-2\lambda/B)\le\bar p\). The source conditional mean at an edge
has local norm at most \(\bar p\), because it reaches distance at most one.
Choose rational upper bounds \(\bar e_\beta\ge e^\beta\),
\(\bar e_{2\lambda}\ge e^{2\lambda}\), and \(\bar e_\kappa\ge e^\kappa\),
and set \(H_B=\bar e_\beta\bar e_{2\lambda}\mathbb A_B\). Each observed
strip has radius at most \(B\) from any of its members, hence local
activity at most \(H_B\).

A fine edge belongs to at most one observed strip. If it does, write
\(\bar z_p=\bar z_p^{(0)}+z_e\bar z_p^{(1)}\) and let \(f_e\) be its
conditional mean under the source. Normalizing the two probabilities after
multiplying by the observed factor gives
\begin{equation}
 F_e=\frac{f_e+y_p(f_e\bar z_p^{(0)}+\bar z_p^{(1)})}
          {1+y_p(\bar z_p^{(0)}+f_e\bar z_p^{(1)})}
     =f_e+\frac{y_p(1-f_e^2)\bar z_p^{(1)}}
               {1+y_p(\bar z_p^{(0)}+f_e\bar z_p^{(1)})}.
 \label{eq:rg-spatial-proof-normalized-tilted-row}
\end{equation}
If no observed strip touches the edge, \(F_e=f_e\); the target strip is
omitted when computing its own conditional mean. Both
\(y_p\bar z_p^{(0)}\) and \(y_p\bar z_p^{(1)}\) have local norm at most
\(H_B\). Increase \(B\) until
\begin{equation}
 H_B\le\min\left\{\frac1{2(1+\bar p)},
                       \frac{1-\bar p}{4(1+\bar p^2)}\right\}.
 \label{eq:rg-spatial-proof-basin-row-gate}
\end{equation}
The denominator perturbation then has norm at most \(1/2\), so its
geometric inverse has norm at most two, and the correction has norm at most
\(2(1+\bar p^2)H_B\le(1-\bar p)/2\). Every normalized conditional mean
therefore has local norm at most \(r=(1+\bar p)/2<1\); the denominator and
its even contribution are retained.

Let \(\Pi\) project onto terms with empty fine support. For nonempty \(S\),
define
\begin{equation}
 \mathcal A(\chi_Sy_T)=\frac1{|S|}\sum_{e\in S}
              \chi_{S\setminus\{e\}}y_TF_e,\qquad
 \mathcal A_\perp=(1-\Pi)\mathcal A,
 \label{eq:rg-spatial-proof-absorbing-operator}
\end{equation}
and let \(\mathcal A\) vanish on fine-constant terms. Removing the fine
edge, using its local row norm, and averaging give
\(\|\mathcal A(1-\Pi)\|\le e^{-\alpha}r\le r\); the average avoids a
support-size factor. Conditional expectation under the normalized finite
tilted law satisfies \(\E\chi_Sy_T=\E\mathcal A(\chi_Sy_T)\), and absorbing
the fine-constant terms once and iterating proves
\begin{equation}
 \E F=\Pi F+\sum_{k\ge1}\Pi\mathcal A\mathcal A_\perp^{k-1}(1-\Pi)F.
 \label{eq:rg-spatial-proof-absorbing-series}
\end{equation}
The unabsorbed remainder is bounded geometrically in the same norm. At any
fixed substitution order only finitely many fine variables and coarse
labels are reached, and the coefficients stabilize once the finite
inventories contain that region. Fixing exterior signs decreases norms, and
finite weights include every crossing original and observed factor, so the
uniform geometric remainder proves boundary-independent continuous limits
for arbitrary full observed patterns.

For a finite target set \(A\) of coarse edges with target signs \(y_A\),
observing its complement gives the conditional density
\(\E_{\mu^{y,A^c}}\prod_{e\in A}(1+y_e\bar z_e)\) relative to fair counting
on \(A\). It is normalized and lies between \((1-\mathbb A_B)^{|A|}\) and
\((1+\mathbb A_B)^{|A|}\). Increasing finite observed subsets gives this
formula uniformly for every exterior by the preceding stabilization and
tail, finite Bayes consistency passes to these limits, and testing their
conditional integrals against exterior cylinders identifies them as the DLR
kernels of the pushed law. This supplies the full finite kernels, not only
one mean or an almost-sure version.

For a single target, start \cref{eq:rg-spatial-proof-absorbing-series} at
\(\bar z_{p_0}\), whose rooted norm is at most
\(\bar e_{2\lambda}\mathbb A_B\). Its conditional mean \(m_{p_0}\)
therefore has norm at most \(M_B=\bar e_{2\lambda}\mathbb A_B/(1-r)\).
Increase \(B\) until \(M_B\le1/2\). The norm-convergent series
\(\widetilde h_{p_0}=\operatorname{atanh}m_{p_0}\) gives the complete
half-log-odds, with \(\|\widetilde h_{p_0}\|\le\bar h_B=M_B/(1-M_B^2)\).
Positive two-edge conditional kernels satisfy the conditional square
identity
\begin{equation}
 \widetilde h_e(y^{f,+})-\widetilde h_e(y^{f,-})
 =\widetilde h_f(y^{e,+})-\widetilde h_f(y^{e,-}),
 \label{eq:rg-spatial-proof-conditional-square}
\end{equation}
since both sides are half the same log cross ratio. Taking Walsh
coefficients, justified by absolute convergence, shows that the coefficient
of \(\chi_{T\setminus\{e\}}\) in \(\widetilde h_e\) does not depend on the
marked edge \(e\in T\); define it to be \(J_T^{(B)}\). It reproduces every
single-edge ratio, and connectedness of the finite sign cube with
normalization then reproduces all the conditional kernels. Translation
covariance gives one canonical orbit representative.

For a size-\(n\) orbit, its translated incidence and norm prefactor are at
most \(n^2\). The inequality \(n^2e^{\kappa n}\le2e^\kappa
e^{(\kappa+1)(n-1)}\) follows from \(\sup_{s\ge1}s^2e^{-(s-1)}=4/e<2\).
Because the root belongs to \(T\), its diameter is at most twice its rooted
radius. Selecting only supports whose literal first edge is the root and
summing the two root orientations therefore gives
\begin{equation}
 \|J^{(B)}\|_{\kappa,\lambda}
 \le\mathcal B_B:=4\bar e_\kappa\frac{M_B}{1-M_B^2}.
 \label{eq:rg-spatial-proof-basin-norm}
\end{equation}
This estimate counts all supports and translated incidences in the
prescribed weights; the fine weight \(\alpha\) was auxiliary, and
\(\lambda\) has not been replaced by a smaller output weight.

Every displayed gate is eventually true and remains true for larger \(B\):
the base gate improves with \(B\), and \(\mathbb A_B\) decreases
exponentially. Rational exponential enclosures as in
\cref{eq:rg-spatial-proof-exp-enclosure} give a positive rational lower
bound \(\varepsilon_0<\varepsilon\), and a finite rational search for
\(B_0\) ensuring the gates and \(\mathcal B_B<\varepsilon_0/16\)
terminates, uniformly over the compact box. The local theorem,
\cref{thm:rg-spatial-local-operator}, now supplies uniqueness in its ball,
and extraction of conditional coefficients with path composition identifies
\(\mathscr R_aJ^{(B)}=J^{(aB)}\).

To obtain a finite rational interaction, keep absorbed substitution orders
through \(n_1\), obtaining \(m^{[n_1]}_{p_0}\), and keep the odd powers of
\(\operatorname{atanh}m^{[n_1]}_{p_0}\) through \(2n_2-1\). For rational
\(b,t\), every local row in \cref{eq:rg-spatial-proof-normalized-tilted-row}
is a rational function of finitely many signs with positive denominator.
Evaluate it on the complete cube and take its finite Walsh transform,
retaining every coefficient. The cube may be reduced exactly to a basis for
the \(\mathbb F_2\)-span of its character supports: the basis-character map
is surjective by row rank, the function is constant on its fibers, and
every outside-span coefficient vanishes by pairing with a kernel element
that flips that character. The transform on the basis cube, lifted to
literal supports, is therefore the entire function of the literal
coordinates; the weights still use literal cardinalities and distances, not
the number of basis generators. The local strip row's character rank is at
most \(2B\), from its \(2(B-2)\) three-edge characters, endpoint character,
possibly removed edge, and two original row characters.

All finite substitution orders and powers have finite inventories. The mean
remainder is \(d_{n_1}=M_Br^{n_1}\), and both exact and truncated absorbed
sums have norm at most \(M_B\). Telescoping the odd powers gives the
\(\operatorname{atanh}\) Lipschitz bound \((1-M_B^2)^{-1}\), and its omitted
powers give the error
\begin{equation}
 \|\widehat J-J^{(B)}\|_{\kappa,\lambda}
 \le4\bar e_\kappa\left[
       \frac{d_{n_1}}{1-M_B^2}
       +\frac{M_B^{2n_2+1}}{(2n_2+1)(1-M_B^2)}\right].
 \label{eq:rg-spatial-proof-basin-oracle-error}
\end{equation}
The incidence argument applies also to omitted coefficients and to
approximate rows that need not obey exact square identities, and canonical
selection gives one valid approximate interaction. After \(n_1\)
substitutions its labels lie within coarse radius \(n_1+1\), and the full
\(\operatorname{atanh}\) of that finite-variable mean introduces no new
labels, so supports of diameter greater than \(2(n_1+1)\) have total norm at
most \(4\bar e_\kappa d_{n_1}/(1-M_B^2)\). The finite odd polynomial pays
its separate power tail in \cref{eq:rg-spatial-proof-basin-oracle-error}.

Increasing \(n_1,n_2\) makes the error arbitrarily small, and each rational
transform, substitution, and product terminates. For a finite word, choose
the input error small enough that both the true interaction and its
approximation lie in one ball of radius \(r<\varepsilon\). Comparing their
exact trajectories by the Lipschitz bounds gives
\(\delta\prod_k\operatorname{Lip}_{a_k}(r)\), and adding the separate Taylor
word error proves the stated total. The argument retains normalized fine
tilts throughout and computes no discarded scalar partition.
\hfill\(\square\)

\section{Higher response and the Fisher-density geometry}
\label{sec:rg-spatial-proof-fisher}

\noindent\emph{Proof of \cref{thm:rg-spatial-fisher-density}.}
Enlarge the compact rectangle of natural parameters to
\(\widehat K=[h_-,h_+]\times[0,j_+]\), with \(h_->0\) and
\(0<j_+<\operatorname{atanh}(1/12)\). Its positive-\(j\) interior contains
the original compact set, and the boundary \(j=0\) is used only for the
one-sided continuation below. Write \(t_+=\tanh j_+\), choose \(\tau>0\)
with \(\mathrm q=12t_+e^\tau<1\), and put \(\mathrm c_1=16\mathrm q/(1-\mathrm q)\),
\(\rho=\tanh(h_++2j_+)\), and \(d=1-\rho>0\). These constants work for
every volume, depth, support pattern, and fixed fine exterior. Finite
supports and references are parameter independent. The finite density is
an exponential family in \((h,j)\); its field statistic sums free edge
signs, and its plaquette statistic includes every touching plaquette with
its exterior completion. Expectations and cumulants under the fine law at
\((h,j)\) carry the subscript \(h,j\).

\subsection{Localized source cumulants and derivatives}
\label{sec:rg-spatial-proof-higher-response}

Let \(F\) be a finite vector of block observables supported on an anchor
set \(A\), and \(F'\) another supported outside its radius-\(R\)
neighborhood. Condition on the fine blocks outside that neighborhood. The
weighted Hamming proof compares any two resulting interior rows. If the
finite whole-block volume is irregular, its inherited exterior is identical
on both sides and contributes zero disagreement; only the newly revealed
shell contributes. Thus
\begin{equation}
 \operatorname{TV}(\operatorname{Law}(F,F'),
       \operatorname{Law}(F)\otimes\operatorname{Law}(F'))
 \le\min\{1,\mathrm c_1|A|B^2(R+1)e^{-\tau(R+1)}\}.
 \label{eq:rg-spatial-proof-group-decoupling}
\end{equation}
Each conditional law of \(F\) is within this bound of their mixture, and
integrating with the exterior and marginalizing proves the group
comparison. It passes to infinite constrained laws by finite-cylinder
convergence, also at \(j=0\) by the exact product-fiber law.

For \(m\) bounded observables define
\begin{equation}
 \operatorname{cum}(F_1,\ldots,F_m)
 =\sum_{\pi\in\mathcal P_m}(-1)^{|\pi|-1}(|\pi|-1)!
                              \prod_{D\in\pi}\E\prod_{i\in D}F_i,
 \label{eq:rg-spatial-proof-response-cumulant}
\end{equation}
where \(\mathcal P_m\) is the set of partitions of \(\{1,\ldots,m\}\). Its
absolute value is at most \(k_m\prod_i\|F_i\|_\infty\), and a change of the
joint law by total variation \(v\) changes it by at most
\(\bar k_mv\prod_i\|F_i\|_\infty\), where \(k_m=\sum_\pi(|\pi|-1)!\) and
\(\bar k_m=2\sum_\pi|\pi|!\). For \(m=2,3,4\), these pairs are
\((2,6),(6,26),(26,150)\). The second estimate telescopes each product of
moments, each moment changing by at most twice the total variation. A
cumulant of two independent nonempty groups vanishes, because its
generating function factorizes and its logarithm has no mixed coefficient;
expanding that finite logarithm also proves
\cref{eq:rg-spatial-proof-response-cumulant}.

Let a local \(f\) have anchor support \(A\) of size \(a_0\ge1\), and let
\(r=1,2,3\) labeled edge or plaquette sources have base anchors
\(x_1,\ldots,x_r\). Put \(M=\max_s\operatorname{dist}_\infty(x_s,A)\). On the
complete graph of these labeled anchors and a root representing \(A\), a
minimum-weight spanning tree has an edge of length at least \(M/r\).
Removing a longest edge splits the graph into groups with all cross
distances at least that length; otherwise exchanging a shorter cross edge
would reduce the tree. Plaquette supports enlarge anchor sets by at most one
on either side. For \(M\ge4r\), use \(R=\lfloor M/r\rfloor-3\) in
\cref{eq:rg-spatial-proof-group-decoupling}; the root group has at most
\(a_0+4r\) anchors. For smaller \(M\), use the individual cumulant bound.
With
\begin{equation}
 L_r(a_0)=e^{4\tau}[k_{r+1}+\bar k_{r+1}\mathrm c_1(a_0+4r)],\qquad
 S_r(a_0)=\frac{(8a_0)^rL_r(a_0)(2r+1)!}
                  {(1-e^{-\tau/r})^{2r+2}},
 \label{eq:rg-spatial-proof-response-constants}
\end{equation}
this proves the bound \(L_r(a_0)B^2(M+1)e^{-\tau M/r}\|f\|_\infty\) for
each cumulant. There are at most \(2B^2\) choices of a source of one
derivative type at each anchor, and at most \(a_0(2M+1)^2\) anchors within
distance \(M\). Counting ordered tuples, including repeats, gives
\begin{equation}
 \begin{split}
 \sum_{s_1,\ldots,s_r}|\operatorname{cum}(f,T_{s_1},\ldots,T_{s_r})|
 &\le S_r(a_0)B^{2r+2}\|f\|_\infty,\\
 \sum_{\max\operatorname{dist}(x_s,A)>L}
       |\operatorname{cum}(f,T_{s_1},\ldots,T_{s_r})|
 &\le S_r(a_0)B^{2r+2}(L+2)^{2r+1}
                     e^{-\tau(L+1)/r}\|f\|_\infty.
 \end{split}
 \label{eq:rg-spatial-proof-response-tails}
\end{equation}
Here \(T_s\) is the edge or plaquette source for its labeled derivative
direction. To sum the bound, use \((M+1)^k\le k!\binom{M+k}k\) in the
geometric series; for the tail, write \(M=L+1+n\) and bound
\((M+1)^k\le(L+2)^k(n+1)^k\). Exterior constant sources have zero cumulant
and may safely be overcounted.

Finite differentiation of the normalized exponential family gives these
cumulant sums for each derivative of a parameter-independent observable.
For an observable \(f_{h,j}\) that depends on the parameters,
\begin{equation}
 \partial^\alpha\E_{h,j}f_{h,j}
 =\sum_{\beta\le\alpha}\binom\alpha\beta
   \sum_{s_1,\ldots,s_{|\alpha-\beta|}}
   \operatorname{cum}_{h,j}(\partial^\beta f_{h,j},T_{s_1},\ldots,
                                      T_{s_{|\alpha-\beta|}}),
 \label{eq:rg-spatial-proof-moving-response}
\end{equation}
where an empty source list means ordinary expectation. Differentiating the
finite generating function or its normalized finite sum proves all
multiplicities, including repeated directions. For each fixed source
tuple, finite-cylinder convergence gives convergence of its moments,
uniformly on \(\widehat K\); near tuples form a finite set, and far tuples
have the uniform tails in \cref{eq:rg-spatial-proof-response-tails}.
Sending first the source cutoff and then the volume to infinity proves
uniform convergence of finite derivatives through order three, and the sums
are continuous by the same finite-near, uniform-tail argument. Passing
finite coordinate fundamental-theorem-of-calculus identities to the limit
proves that these are continuous derivatives, successively through order
three, with agreeing mixed derivatives. At \(j=0\) the argument runs on
right-sided coordinate segments and gives one-sided \(C^3\) continuation.
No infinite normalizer or null-event probability is differentiated.

\subsection{Full local jets and interaction derivatives}
\label{sec:rg-spatial-proof-local-jets}

For a free path, its conditional mean \(\bar z\) given the other fine blocks
is the product of \(B\) hyperbolic tangents with fields \(h+jw_e\),
\(|w_e|\le2\), on at most nine anchor cells. For \(n\le3\),
differentiating the product in labeled directions gives
\begin{equation}
 \|\partial^\beta\bar z\|_\infty
 \le g_nB^n\rho^{B-n},\qquad g_n=2^nn!,\quad |\beta|=n.
 \label{eq:rg-spatial-proof-product-jets}
\end{equation}
The first three hyperbolic-tangent derivatives have bounds \(1,2,6\), and
each \(j\) differentiation adds a factor at most two. There are at most
\(B^n\) assignments to product factors, and the product of their factorials
is at most \(n!\). Negative exponents when \(B<n\) only enlarge the upper
bound. Apply
\cref{eq:rg-spatial-proof-moving-response,eq:rg-spatial-proof-response-tails}
to \(\bar z\). With
\begin{equation}
 D_n=g_n+\sum_{b'=0}^{n-1}\binom n{b'}g_{b'}S_{n-b'}(9),\qquad
 M_n(B)=D_nB^{2n+2}\rho^{B-n},\quad M_0(B)=\rho^B,
 \label{eq:rg-spatial-proof-mean-jets}
\end{equation}
every order-\(n\) derivative of a free-path mean is bounded by \(M_n(B)\).
These bounds tend to zero as \(B\) increases.

For patterns agreeing through \(Q_R\), the cell square of radius \(R\),
choose \(a_R=\lfloor R/2\rfloor-1\) when \(R\ge2\). A near source tuple has
joint support within \(Q_{a_R+1}\); apply the group comparison on its
remaining buffer and count \([2B^2(2a_R+1)^2]^n\) ordered tuples, and use
\cref{eq:rg-spatial-proof-response-tails} for each of the two far sums.
Writing \(\mathrm c'_n=2^n\bar k_{n+1}\mathrm c_1(1+4n)\), the resulting derivative modulus
is
\begin{equation}
\begin{split}
 V_n(B,R)=\min\bigl\{2M_n(B),B^{2n+2}\bigl[
 &\mathrm c'_n(2a_R+1)^{2n}(R-a_R)e^{-\tau(R-a_R)}\\
 &+2S_n(1)(a_R+2)^{2n+1}e^{-\tau(a_R+1)/n}\bigr]\bigr\}.
\end{split}
 \label{eq:rg-spatial-proof-jet-modulus}
\end{equation}
For \(R=0,1\) take \(V_n=2M_n\), and for values use
\(V_0=\min\{2\rho^B,2\mathrm c_1B^2(R+1)e^{-\tau(R+1)}\}\). The same proof
compares a finite whole-block conditional calculation against its infinite
kernel: near sources are interior, and finite crossing sources and infinite
far sources satisfy the two separate tail bounds. For \(R\ge2\), the
bracket in \cref{eq:rg-spatial-proof-jet-modulus} is at most
\(Z_n(R+1)^{2n+1}e^{-\tau R/(2n)}\), with
\(Z_n=\mathrm c'_n+2e^{\tau/(2n)}S_n(1)\).

Every Bernoulli factor is \((1+s\bar z)/2\), with \(1+s\bar z\ge d\).
Revealing both target parities sequentially retains all their dependence,
and their positive four-state row is \(C^3\). For one Bernoulli log factor
the derivative amplitudes are
\begin{equation}
 A_1=\frac{M_1}{d},\qquad
 A_2=\frac{M_2}{d}+\frac{M_1^2}{d^2},\qquad
 A_3=\frac{M_3}{d}+\frac{3M_2M_1}{d^2}+\frac{2M_1^3}{d^3}.
 \label{eq:rg-spatial-proof-log-jet-amplitudes}
\end{equation}
The multiplicities follow by differentiating \(\log(1+s\bar z)\) three
times. Its derivative moduli are
\begin{equation}
 \begin{split}
 T_1={}&V_1/d+M_1V_0/d^2,\\
 T_2={}&V_2/d+(M_2V_0+2M_1V_1)/d^2+2M_1^2V_0/d^3,\\
 T_3={}&V_3/d+(M_3V_0+3M_2V_1+3M_1V_2)/d^2\\
       &+(6M_2M_1V_0+6M_1^2V_1)/d^3+6M_1^3V_0/d^4.
 \end{split}
 \label{eq:rg-spatial-proof-log-jet-moduli}
\end{equation}
Subtracting the derivative ratios and using \(|r_1^{-k}-r_2^{-k}|\le
k|r_1-r_2|/d^{k+1}\) for \(r_1,r_2\ge d\) gives these formulas. One-cell log
probabilities use \(2A_n,2T_n\), and reference-state log odds use
\(4A_n,4T_n\). Each modulus can also be capped by twice its amplitude. The
same already revealed first parity is fixed when comparing the second
factor, so the support patterns still agree in their common interior.

The complete interaction jets are summable, with explicit constants. Set
\(E_1=D_1/d\), \(E_2=D_2/d+D_1^2/d^2\), and
\(E_3=D_3/d+3D_2D_1/d^2+2D_1^3/d^3\). Then
\(F_n(B)=4A_n(B)\le4E_nB^{4n}\rho^{B-n}\). Each \(M_j\) is bounded over all
\(B\) by \(U_j=D_j\rho^{1-j}(2j+2)!/(1-\rho)^{2j+3}\). Put
\(Z_0=2\mathrm c_1\), and obtain \(P_n\) from the right side of
\cref{eq:rg-spatial-proof-log-jet-moduli} by replacing every \(M_j\) with
\(U_j\) and every \(V_j\) with \(Z_j\); these are explicit positive
constants, and the powers of \(d\) are unchanged. For \(R\ge2\),
\(T_n\le P_nB^{2n+2}(R+1)^{2n+1}e^{-\tau R/(2n)}\). The constant
\begin{equation}
 H_n=4P_n+
 8E_n\rho^{-n}\frac{(2n-2)!}{(1-\rho)^{2n-1}}e^{\tau/(2n)}
 \label{eq:rg-spatial-proof-jet-shell-constant}
\end{equation}
also covers shell radii one and two through the amplitude cap. Its
order-\(n\) canonical shell is therefore bounded by both \(2F_n(B)\) and
\(H_nB^{2n+2}r^{2n+1}e^{-\tau(r-1)/(2n)}\). Choose one common
\(\lambda_f=\tau/48\). Geometric interpolation,
\(\operatorname{diam}A_{x,r}\le2r\), and the translate count give
\begin{equation}
 \|\partial^\alpha\mathscr U_B\|_{\lambda_f}
 \le F_n(B)+\sqrt{2F_n(B)H_n}B^{n+1}e^{\tau/(4n)}
       \sum_{r\ge1}(2r+1)^2r^{n+1}
          e^{-[\tau/(4n)-2\lambda_f]r},\quad |\alpha|=n\le3.
 \label{eq:rg-spatial-proof-complete-jet-norm}
\end{equation}
The bracketed exponent \(\tau/(4n)-2\lambda_f\) is positive, the positive
series is at most \(9(n+3)!\) divided by
\((1-e^{-[\tau/(4n)-2\lambda_f]})^{n+4}\), and its omitted tail has an
explicit convergent bound. The right side decays as a constant times
\(B^{3n+1}\rho^{(B-n)/2}+B^{4n}\rho^{B-n}\). Values use the canonical
interaction proof of \cref{sec:rg-spatial-proof-potential} in the same
weaker \(\lambda_f\) norm. Finite-shell maps are \(C^3\), and their
derivatives through order three converge uniformly in this Banach norm, so
successive coordinate fundamental-theorem-of-calculus identities prove that
the entire interaction is \(C^3\). No shell or derivative is replaced by a
finite truncation in the statement.

\subsection{Moving expectations and marginal information}
\label{sec:rg-spatial-proof-moving-fisher}

For any finite or infinite subset \(S\) of the lexicographic past, let
\(p_S(s;y)\) be the four-state predictive row of the target state \(s\) with
exactly the past coordinates \(y_S\) prescribed and the other paths,
including the future, free. Let \(l=\log p_S\), with subscripts denoting
derivatives in the natural parameters, and use the indices \(i,m,k,n\) for
those derivatives. Define \(\mathsf I_{S,im}=\sum_sp_Sl_il_m\), an
observable of the prescribed past. Its first and second parameter
derivatives, keeping \(y\) fixed, are
\begin{equation}
 \begin{split}
 \mathsf I_{S,im,k}
 &=\sum_sp_S(l_kl_il_m+l_{ik}l_m+l_il_{mk}),\\
 \mathsf I_{S,im,kn}
 &=\sum_sp_S\bigl[(l_{kn}+l_kl_n)l_il_m
 +l_k(l_{in}l_m+l_il_{mn})+l_n(l_{ik}l_m+l_il_{mk})\\
 &+l_{ikn}l_m+l_{ik}l_{mn}+l_{in}l_{mk}+l_il_{mkn}\bigr].
 \end{split}
 \label{eq:rg-spatial-proof-fisher-observable-jets}
\end{equation}
These formulas differentiate the probabilities as well as the scores. Put
\(L_i=2A_i\) for \(i=1,2,3\),
\begin{equation}
 E_B^F=\max\{L_1^2,L_1^3+2L_1L_2,
                L_1^4+5L_1^2L_2+2L_1L_3+2L_2^2\},\qquad
 H_B^F=\max(1,L_1,L_2,L_3).
 \label{eq:rg-spatial-proof-fisher-observable-bound}
\end{equation}
Every matrix entry and its first two derivatives are bounded by \(E_B^F\).
For patterns agreeing in \(Q_R\), every row probability changes by at most
\(V_0\), and each log jet by at most \(2T_n\). Telescoping the at most ten
expanded monomials in \cref{eq:rg-spatial-proof-fisher-observable-jets},
with at most five factors per term and four states, gives the common
modulus
\begin{equation}
 \omega_B(R)=\min\{2E_B^F,
 200(H_B^F)^4[V_0(B,R)+2T_1(B,R)+2T_2(B,R)+2T_3(B,R)]\}.
 \label{eq:rg-spatial-proof-fisher-observable-modulus}
\end{equation}
At fixed \(B\) it is bounded by an explicit constant times
\((R+1)^7e^{-\tau R/6}\). These are differences of jets, not derivatives of
an earlier scalar continuity bound.

Expectations now use the original unconstrained fine law, with coarse
observables pulled back by the fixed path map. For a local \(C^2\)
observable on \(Q_N\), put \(a_N=(2N+1)^2\) and
\(W_B(N)=1+2S_1(a_N)B^4+S_2(a_N)B^6\). The moving response formula gives
\begin{equation}
 \max_{|\alpha|\le2}\sup_{(h,j)\in\widehat K}
      |\partial^\alpha\E_{h,j}f_N|
 \le W_B(N)\max_{|\alpha|\le2}\sup_{(h,j),y}|\partial^\alpha f_N|.
 \label{eq:rg-spatial-proof-varying-expectation-bound}
\end{equation}
In particular, the second derivative includes \(\E f_{im}\), both source
covariance sums involving \(f_i,f_m\), and the complete three-variable
cumulant sum involving \(f\). Formula
\cref{eq:rg-spatial-proof-response-tails} justifies each passage from finite
volumes, and neither the observable nor the law is held fixed incorrectly.
The function \(W_B(N)\) is polynomial in \(N\). For a shell series
\(f=\sum_Nd_N\) with \(\sum_NW_B(N)\|d_N\|_{C^2}<\infty\), finite partial
sums and their expectation derivatives are uniformly Cauchy, and repeated
fundamental-theorem-of-calculus identities give \(C^2\) of its expectation
with termwise differentiation.

Apply this to \(f_N^S=\mathsf I_{S\cap Q_N}\). Its initial shell is bounded
by \(E_B^F\), and its later shell differences through order two by
\(\omega_B(N-1)\). Thus
\begin{equation}
 Z_B^F=W_B(0)E_B^F+\sum_{N\ge1}W_B(N)\omega_B(N-1)<\infty,
 \qquad
 \epsilon_B(R)=2\sum_{N>R}W_B(N)\omega_B(N-1)\longrightarrow0.
 \label{eq:rg-spatial-proof-fisher-expectation-tails}
\end{equation}
If two past sets agree inside \(Q_R\), their finite prefixes agree exactly
through that radius, so their full-law Fisher-observable expectations
differ in entrywise \(C^2\) norm by at most \(\epsilon_B(R)\), uniformly
over the remote shape of the past. In particular the full-past expectation
\(\mathcal I_B=\E\mathsf I_{\mathrm{past}}\) is \(C^2\).

The positive marginal on \(\Lambda_L\) factors into its predictive rows in
lexicographic order. Each score increment is centered conditionally on its
earlier revealed cells, and two increments are orthogonal because the
earlier one is measurable for the later past. The finite Fisher matrix is
therefore the sum of expected row-Fisher matrices with their finite pasts,
each a \(C^2\) function by the preceding varying-expectation argument.
Translation invariance removes the shifts. For \(L\ge2R\), all but at most
\(4RL\) cells have radius \(R\) inside the square, and their finite pasts
and full past agree on that radius; each remaining difference is bounded by
\(2Z_B^F\). Therefore
\begin{equation}
 \|L^{-2}I_{B,\Lambda_L}-\mathcal I_B\|_{C^2,\mathrm{entry}}
 \le\epsilon_B(R)+8Z_B^FR/L.
 \label{eq:rg-spatial-proof-actual-fisher-limit}
\end{equation}
Sending first \(L\) to infinity at fixed \(R\), then \(R\) to infinity,
proves the marginal \(C^2\) limit, including all boundary-score
derivatives. The canonical row with future fixed to vacuum is never
substituted for the predictive row.

\subsection{An exhibited positive chart and its curvature}
\label{sec:rg-spatial-proof-regular-chart}

At \(j=0\), with any fixed complete past, the blocks are independent
product parity fibers, and the two target path parities are independent
with mean \(b^B\). For one fine edge of a target path with parity \(Y\),
\begin{equation}
 a(s)=\E[z_e\mid Y=s]-b
      =\frac{sb^{B-1}(1-b^2)}{1+sb^B},\qquad
 v_B=\E a^2=\frac{b^{2B-2}(1-b^2)^2}{1-b^{2B}}>0.
 \label{eq:rg-spatial-proof-product-score-contrast}
\end{equation}
Multiplying the product law by its parity indicator and summing proves this
formula; \(\E a=0\), and conditional-variance decomposition gives
\(v_B\le1-b^2\le1\). The field score is \(\mathbf s_h=B(a_1+a_2)\), and the
plaquette score is \(\mathbf s_j=b^2a_1a_2+\alpha_1a_1+\alpha_2a_2\), with
\(|\alpha_i|\le2B\). For the mixed term, the plaquette at the common
starting vertex is the unique one meeting both outgoing target paths; its
other two edges are future paths at \(B=1\), or unused singleton edges at
\(B>1\), and have mean \(b\), so its centered contribution is
\(b^2a_1a_2+b^3(a_1+a_2)\). Every other incident plaquette meets only one
target path and contributes a multiple of \(a_i\) of magnitude at most one;
there are \(2B-1\) such incidences. Nonincident plaquettes contribute zero
by product-block independence. These are derivatives of the full family
through its source formula and one-sided limit, not derivatives of a
formula restricted to \(j=0\).

Independence and centering make \(a_1a_2\) orthogonal to both linear
contrasts. The conditional Fisher entries are therefore
\(I_{hh}=2B^2v_B\), \(I_{hj}=Bv_B(\alpha_1+\alpha_2)\), and
\(I_{jj}=v_B(\alpha_1^2+\alpha_2^2)+b^4v_B^2\). Their determinant is
\(B^2v_B^2(\alpha_1-\alpha_2)^2+2B^2b^4v_B^3\), and their trace is at most
\(11B^2v_B\). Dividing the determinant by the trace bounds the smaller
eigenvalue below by \(2b^4v_B^2/11\), uniformly in the past. On a compact
positive-field interval with \(b\in[b_-,b_+]\), this is at least
\begin{equation}
 \eta_{0,B}=\frac2{11}b_-^{4B}(1-b_+^2)^4>0.
 \label{eq:rg-spatial-proof-fisher-product-floor}
\end{equation}

Use the bounds \(A_1,A_2\) on the enlarged rectangle and put
\(Q_B=16(A_1^3+A_1A_2)>0\). The score norm is at most \(2\sqrt2A_1\), and
its \(j\) derivative at most \(2\sqrt2A_2\). Differentiating the four-state
Fisher sum gives an operator-norm derivative at most \(Q_B\), including the
probability derivative term. One-sided \(C^3\) and the fundamental theorem
of calculus imply
\begin{equation}
 \mathcal I_B(h,j)\succeq\frac{\eta_{0,B}}2I_2
 \quad\text{for}\quad
 0<j<\delta_B:=\min\{j_+/2,\eta_{0,B}/(2Q_B)\}.
 \label{eq:rg-spatial-proof-regular-fisher-region}
\end{equation}
The same floor holds rowwise and hence after the full-law expectation. An
enlarged field interval supplies an open neighborhood of each original
compact interval. For a prescribed finite list of depths, take the finite
minima of positive widths and floors. No positive minimum over every integer
depth is asserted.

On a compact subchart of this region write \(g_{mn}=\mathcal I_{B,mn}\), and
use the indices \(k,m,n,r,s\) for the natural parameters. The coefficients
and curvature
\begin{equation}
 \Gamma^r_{nk}=\frac12g^{rs}
        (\partial_ng_{sk}+\partial_kg_{sn}-\partial_sg_{nk}),\qquad
 R^r{}_{kmn}=\partial_m\Gamma^r_{nk}-\partial_n\Gamma^r_{mk}
             +\Gamma^r_{ms}\Gamma^s_{nk}-\Gamma^r_{ns}\Gamma^s_{mk}
 \label{eq:rg-spatial-proof-fisher-connection-curvature}
\end{equation}
define the Levi--Civita connection and its curvature: substitution shows
torsion freedom and metric compatibility, and the cyclic combination of the
compatibility equations recovers \(\Gamma\). Since the metric is \(C^2\)
with a positive floor, its inverse is \(C^2\), the connection is \(C^1\),
and the curvature and its Ricci and scalar contractions are continuous. By
\cref{eq:rg-spatial-proof-actual-fisher-limit}, the metrics
\(L^{-2}I_{B,\Lambda_L}\) converge to \(\mathcal I_B\) in \(C^2\) on a
neighborhood of the subchart, where \(\mathcal I_B\) is positive definite by
\cref{eq:rg-spatial-proof-regular-fisher-region}. These are the hypotheses
of \cref{prop:rg-thermo-label-regular-geometry}, which gives \(C^2\)
convergence of the inverses, \(C^1\) convergence of the connections, and
\(C^0\) convergence of the curvatures, at this fixed \(B\).

These are also local law charts. At a chart point choose a finite marginal
with positive Fisher matrix. Its probability Jacobian has rank two because
\(I=(Dp)^\top\operatorname{diag}(1/p)Dp\), so the map \(P\) from \((h,j)\)
to two suitably chosen atom probabilities has a nonsingular Jacobian at the
point. By the inverse function theorem \citep[Thm.~9.24]{Rudin1976}, \(P\)
has a \(C^1\) local inverse, and the formula
\(D(P^{-1})=(DP\circ P^{-1})^{-1}\) shows inductively that the inverse is
\(C^3\), because \(P\) is. This exhibits a chart of the law without
assuming a global quotient manifold. \hfill\(\square\)
